% Generated draft: requires mathematical and source review.
\documentclass[11pt]{article}
\usepackage[margin=25mm]{geometry}
\usepackage{fontspec}
\setmainfont{Latin Modern Roman}
\usepackage{amsmath,amssymb,mathtools,unicode-math}
\setmathfont{Latin Modern Math}
\usepackage{enumitem,needspace,xurl,hyperref}
\setlength{\parindent}{0pt}
\setlength{\parskip}{5pt}
\setlength{\emergencystretch}{3em}
\newcommand{\N}{\mathbb{N}}
\newcommand{\Z}{\mathbb{Z}}
\newcommand{\Q}{\mathbb{Q}}
\newcommand{\R}{\mathbb{R}}
\newcommand{\C}{\mathbb{C}}
\newcommand{\F}{\mathbb{F}}
\newcommand{\A}{\mathbb{A}}
\newcommand{\PP}{\mathbb{P}}
\newcommand{\E}{\mathbb{E}}
\newcommand{\eps}{\varepsilon}
\newcommand{\dd}{\,\mathrm{d}}
\begin{document}
\section*{UnsolvedMath ID: 800012}
\textbf{Source problem code:} catalogue entry
\subsection{Definitions and mathematical statement}
Under the usual convention, an \emph{ordered semigroup} is a triple \((S,\cdot,\leq)\), where \(S\) is nonempty, multiplication is associative, and \(\leq\) is a partial order compatible with multiplication: for all \(a,b,c\in S\),
\[
a\leq b\ \Longrightarrow\ ca\leq cb\ \text{ and }\ ac\leq bc.
\]
The \emph{trivial order} is understood here to mean the equality order, namely
\[
\forall a,b\in S,\qquad a\leq b\ \Longleftrightarrow\ a=b.
\]
A standard interpretation of a \emph{variety of ordered semigroups} is a class consisting of all ordered semigroups satisfying a specified set of inequalities between nonempty semigroup words, with each inequality required to hold under every assignment of its variables.

\textbf{Problem (V. Makaridina).} Describe all varieties \(\mathcal V\) of ordered semigroups such that
\[
\forall (S,\cdot,\leq)\in\mathcal V\;\forall a,b\in S,
\qquad a\leq b\ \Longrightarrow\ a=b.
\]
Thus the requirement concerns the designated order of every member of \(\mathcal V\).

\textbf{Unresolved qualifications.} The supplied record does not define ordered semigroups, trivial order, or varieties of ordered semigroups. The conventions above make the target precise conditionally; their agreement with the source's intended conventions requires review. In particular, the wording does not explicitly resolve whether the target instead concerns varieties of underlying semigroups for which every compatible partial order is equality. The source adds that some examples are given in reference [Mak4], whose bibliographic details are not supplied.
\subsection{Short English statement}
V. Makaridina asks for a description of all varieties of ordered semigroups whose members always carry the equality order. The definitions of order and variety require confirmation from the source. Examples are attributed to [Mak4], but that reference is not identified in the supplied record.
\subsection{Sources}
\begin{itemize}
\item Local UnsolvedMath registry, permanent ID \(800012\), exact source problem code \texttt{catalogue entry}, title \emph{The Lyapin's Notebook --- Problem 12}; supplied upstream revision \texttt{d371ee917ac9647fbb06cd6fc529deec6f3dda82}.
\item Bershadsky, Kublanovsky, Mashevitzky, \emph{The Lyapin's Notebook: Unsolved Problems in Semigroup Theory} (2026), Problem 12, \url{https://arxiv.org/abs/2604.04763}. The supplied statement attributes the problem to V. Makaridina. The registry's proposer field names the three collectors and gives the year 2026; it does not establish the original proposal date.
\item The supplied provenance records access on 2026-10-16 and extraction by \texttt{arxiv\_latex}. These are supplied provenance details, not a claim of independent source access or verification.
\item Reference [Mak4]: the supplied statement says that some examples are given there. No author, title, publication details, or link for [Mak4] is supplied.
\item The registry labels this problem \emph{open}. This is a supplied status label, not an independently verified assessment of the literature.
\end{itemize}
\end{document}
